\section{Clocked action reconstruction}
\label{sec:v78-introduction}

A clocked endpoint law contains both dynamical information and the
selection bias of its source.  We study the recovery of an absolute
generating action from several such laws when neither their normalizing
constants nor their common positive sampling weight is known.  The
same problem arises for finite reflected trajectories and for return
maps with positive action roofs.  The basic observation has the form
\[
                f_j(z)\ \propto\ a(z)(T_j-W(z)),
                  \qquad T_1<\cdots<T_J,
\]
where $W<T_1$ is nonconstant.  Three clocks recover $W$ exactly.
This exact recovery alone makes the statistical observation equivalent
to action data, after quotienting common reweighting.  The questions
addressed here are what remains stable when the weights are not quite
shared, what can be validated from records, and when a fixed clock
experiment recovers an unknown generating system rather than an action
on an object-dependent gate.

The first main result is a split description of the entire positive
observation quotient, including observations outside the exact model.
Anchored log-density ratios identify this quotient with a Banach space.
On every nonconstant two-anchor chart, an explicit smooth coordinate
change identifies it with an action and a complementary residual
(Theorem~\ref{thm:v79-split}).  The residual vanishes exactly on the
shared-weight model.  Approximate sharing produces a linearly controlled
action error without spatial derivative loss.  We characterize all
clock-dependent perturbations that are exactly confounded with another
action, and identify the tangent and normal components of arbitrary
small perturbations (Theorem~\ref{thm:v79-nuisance} and
Proposition~\ref{prop:v79-tangent}).

These coordinates have statistical consequences beyond exact
identification.  They give a finite-record model-validation test,
including a functional residual already for three clocks.  On bounded
H\"older classes of the weighted-action experiment, with $n$ retained
samples at each fixed clock and clock-dependent log-weight errors of
$C^m$ size at most $\epsilon$, the minimax expected $C^m$ action error is
\[
           \asymp\left(\frac{\log n}{n}\right)^{(\beta-m)/(2\beta+d)}
                        +\epsilon,
                 \qquad \beta>m,
\]
under the explicit nondegeneracy and interior-gate conditions of
Theorem~\ref{thm:v79-minimax}.  The two lower bounds distinguish sampling
error from exact action confounding.  This is an action-risk theorem;
it is not asserted to give the optimal statistical rate for billiard
geometry.

The second main result gives a class-wide geometric clock experiment.
For an arbitrary unknown two-variable function $L$ with uniformly positive
Hessian and uniformly negative mixed derivative, six deadline-conditioned
endpoint laws of its contact suspension determine $L$ on a prescribed
compact square (Theorem~\ref{thm:v79-six-clock}).  The experiment uses
return counts $n$ and $n+1$, for any prescribed finite $n$, fixed gates,
fixed deadlines and a fixed independent external release-delay interval.
All settings are chosen from class bounds, not from the unknown function.
An exact contact-suspension construction makes the relevant action a
return-time roof; the observer counts returns and compares a physical
clock with a deadline, without evaluating an unknown action.  Determinism,
positive stationary Hessians, twist, uniform gates and matching margins
are proved from inequalities on $L$, rather than included as abstract
reconstruction assumptions.  This class contains nonseparable and
nonanalytic smooth perturbations.  It also contains the retained
mass--potential example, which recovers the potential on a prescribed
compact interval.

The third stage is the geometric inverse for marked dispersing billiards.
Their physical time windows yield the same clock laws.  Recovering two
successive branch actions, matching their initial covectors, and
subtracting the common prefix gives a final-flight distance kernel.
A finite distance certificate reconstructs the boundary arcs, and
labelled overlaps register them in one Euclidean plane.  Whole-table
coverage follows from a finite identifying atlas.  This geometric
conclusion, including nonanalytic smooth boundaries, is retained in
full below.  Its uniform estimates concern bounded neighborhoods of
a fixed reference atlas, not a single protocol for all billiards.
The final theorem adds clock-dependent misspecification to its
finite-record error bound and explicitly charges every raw launch,
including the cost of branch classification.

\subsection{The marked physical experiment}
A periodic table is the complement of disjoint translates
$K_a+L\lambda$, $1\le a\le q$, $\lambda\in\Z^2$, where $L$ is
invertible and each obstacle has a smooth strictly convex boundary with
positive curvature.  A scalar boundary atlas, obstacle and lift labels,
and regular reflection words are part of the observation convention.
An atlas coordinate is realized by an unknown regular embedding
$R_I:I\to\R^2$.  Its Euclidean position, speed, tangent and normal are
not observations.  The speed of the moving particle is one.  One can
instead fix, once for the whole class, the abstract label space
$\{1,\ldots,q\}\times(\R/\Z)\times\Z^2$ and its standard charts;
then chart transitions are fixed by this label space, not chosen from
an unknown Euclidean realization.

A branch $\gamma$ has first and last scalar labels $(s,t)$ and
stationary length $W_\gamma(s,t)$.  The observer knows its word and its
endpoint gate, but records only $(s,t)$ on retention.  An unsuccessful
attempt has a failure symbol.  In the flow box before the first collision,
the source has an unknown smooth positive density $\rho_\gamma(s,p)$,
independent of the time $\alpha$ to that collision.  The unknown
retention probability $e_\gamma(s,t)>0$ is common to the clocks of this
branch.  It need not separate in $s,t$, and different branches may have
different preparations and retention laws.  All three conditional laws
are normalized on the same gate.  Their restrictions to a smaller gate
are not renormalized.

Let $\tau^-_\gamma,\tau^+_\gamma$ be the adjacent free flights.  Three
absolute times $T_1<T_2<T_3$ are admissible on a gate if, with strict
margins,
\begin{equation}\label{eq:v77-window}
 0<T_j-W_\gamma(z)<\min(\tau^-_\gamma(z),\tau^+_\gamma(z)).
\end{equation}
This is a restriction on the physical class.  It is not a measured
length function.  The resulting densities are
\begin{equation}\label{eq:v78-intro-law}
 f_{\gamma,j}(z)=
 \frac{a_\gamma(z)(T_j-W_\gamma(z))}
 {\int_{Q_\gamma}a_\gamma(u)(T_j-W_\gamma(u))\,du},\qquad
 a_\gamma=e_\gamma\rho_\gamma(s,-W_{\gamma,s})
                   |W_{\gamma,st}|.
\end{equation}
The flow-box proof is given in Lemma~\ref{lem:v77-physical}.
For raw-launch acquisition with classification, the clock-independent
classification probability is an additional factor of $a_\gamma$;
Definition~\ref{def:v79-raw} and Proposition~\ref{prop:v79-raw-cost}
specify its meaning and cost without renormalizing away branch rarity.

A boundary is \emph{locally nonconic} when no nonempty open subarc
is contained in the zero set of a nonzero polynomial of degree at most two.
This is a sufficient condition for the finite certificate, not a
necessity claim for rigidity.

\begin{definition}[Geometric certificate]\label{def:v77-atlas}
A finite identifying atlas consists of finitely many regular prefixes
and their one-flight extensions, fixed scalar gates and admissible clock
triples.  The final-flight patches obtained by matching their common
initial states cover the target.  Each patch has the rank-two and
rank-five distance certificate of Proposition~\ref{prop:v77-distance}.
The graph registering the patches by scalar overlaps is connected.
For a whole periodic table, the cover includes complete quotient
boundaries and common intervals on one boundary and its two marked basis
translates.  Positive matching margins, certificate anchors and overlap
triples are witnesses to this definition, not additional observed marks.
The local construction in Proposition~\ref{prop:v77-cancellation} supplies
such witnesses.  Theorem~\ref{thm:v78-root-free} recovers matching roots
without receiving them.
\end{definition}

\begin{theorem}[Marked statistical-action rigidity]
\label{thm:v78-main}\label{thm:v77-main}
For a finite identifying atlas, the branch-resolved three-clock endpoint
laws determine every target boundary embedding up to a common Euclidean
isometry.  No root-selection bracket, Euclidean contact point, tangent
frame, chart speed, absolute branch length, source weight or lattice
Gram matrix is supplied.  In the whole-table case they determine the
entire marked periodic table and its lattice Gram matrix.

The laws determine their absolute generating actions and the normalized
nuisance shapes separately.  Modulo simultaneous positive reweighting
at the three clocks, their exact information is precisely that of the
absolute actions.  Matching of a prefix to its one-flight extension is
unique on a connected common initial interval whenever it exists.
The matching domain itself is determined by signs of recovered functions.

Every finite-horizon periodic dispersing table with locally nonconic
boundaries admits such a finite protocol, with prefix lengths at least
any prescribed finite $J$.  A single protocol works on a neighborhood
of each such table.  On bounded neighborhoods with positive certificate
margins it has a uniform finite-order stable inverse.  There is also an
explicit finite-record reconstruction by smoothing, root solving,
linear algebra and patch blending; it does not select from an unknown
dense enumeration of possible tables.
\end{theorem}

The theorem compares any two realizations of the same admissible data;
its uniqueness assertion is not a local inverse-function assertion in
the unknown geometry.  Uniform budgets, in contrast, are asserted on the
bounded neighborhoods specified in Section~\ref{sec:v78-uniform}.
The existence of a finite protocol at each table is not a universal
acquisition rule for all unknown tables.  Exact word and lift labels,
and the cross-clock source convention, remain hypotheses throughout.

\subsection{Relation to other inverse data and methods}

For open dispersing billiards, B\'alint--De Simoi--Kaloshin--Leguil
\cite{BalintDeSimoiKaloshinLeguil2020} recover curvature at period-two
points and Lyapunov exponents of periodic trajectories from the marked
length spectrum.  Their observation consists of marked periodic lengths,
not endpoint laws on two-dimensional gates.  De Simoi--Kaloshin--Leguil
\cite{DeSimoiKaloshinLeguil2023} subsequently prove geometric determination
for analytic non-eclipsing open billiards under suitable symmetry and
genericity assumptions.  Our whole-table result concerns smooth
nonanalytic boundaries as well, but uses the different, richer scalar
branch and overlap design stated here.  It is not a theorem replacing
their spectral hypotheses by an intrinsically weaker observation.

Noakes--Stoyanov \cite{NoakesStoyanov2020} prove planar obstacle rigidity
from exterior travelling-time and scattering-length data.  Finamore--Leguil
\cite{FinamoreLeguil} prove finite-horizon Sinai rigidity from an enriched
marked length spectrum, using a CAT(0) comparison construction.  The
external and periodic marks in these results are not the branch-resolved
scalar labels of the present experiment.  The exact comparison actually
proved here is local: a marked canonical graph plus one absolute action
value determines the clock-law orbit, and conversely
(Proposition~\ref{prop:v78-lens}).  The value is essential because a
canonical relation determines derivatives of the action, not its
additive constant.  We make no unproved global ordering of these
observation spaces.

The first-variation interpretation of a discrete Lagrangian, symplecticity
of its map, and stationary composition belong to discrete variational
mechanics; see Marsden--West \cite{MarsdenWest2001}.  In the billiard
setting, the variational relation between lengths and collision
covectors also underlies the marked-length analysis in
\cite{BalintDeSimoiKaloshinLeguil2020}.  We give the finite Hessian and
corner-cofactor calculations needed here, rather than treating these
classical generating identities as new.  The root-selection argument
uses more than a twist inequality at one selected solution: equality of
initial covectors must identify actual trajectories at every zero.  This
physical single-branch condition is stated explicitly in
Theorem~\ref{thm:v78-root-free}.  For the full class
\eqref{eq:v79-convex-class}, it follows globally from convexity and
negative twist.

Positive symplectic actions can be realized as return times in contact
geometry.  Hutchings \cite[Proposition 2.1]{Hutchings2016} constructs an
appropriate contact form from a positive action function of a disk map.
The elementary suspension used here is written out to fix its section,
clock orientation and roof constant.  It is classical infrastructure,
not a new contact-geometric existence theorem.  The inverse application
is the recovery of an unknown two-variable $L$ from one class-wide
six-deadline protocol.  The clock is elapsed Reeb time of that specified
suspension; we do not relabel an arbitrary discrete mechanical action as
the ordinary physical time of the unsuspended map.

The final Euclidean step uses complete bipartite distance information.
Connelly--Gortler--Theran \cite{V77ConnellyGortlerTheran} give a general
theory of generic global rigidity for complete bipartite frameworks.
Our rank-two and rank-five certificate is an explicit sufficient
realization rule, followed by functional trilateration and overlap
registration.  It does not establish necessity of that certificate,
and its failure on some conic configurations is not nonidentifiability.
Its role is to expose the conditioning of a complete smooth-patch inverse.

The nonparametric tools are also standard.  Uniform density estimates
are developed by Gin\'e--Guillou \cite{GineGuillou2002} and
Gin\'e--Nickl \cite{GineNickl2009}, and statistical lower-bound methods
are treated in \cite{Tsybakov2009}.  We supply the kernel and packing
arguments for the precise norm and normalized multi-clock model used
here.  The new structural statement is the explicit split of observations
into action and residual, together with the exact nuisance ambiguity
and its sharp linear contribution to risk.  Neither a quotient of exact
laws nor successful model fitting is a finite-sample sufficiency theorem.

\subsection{Organization and the periodic companion}
Sections~\ref{sec:v77-clocks}--\ref{sec:v79-statistics} develop the
physical clock law, exact information, split quotient and statistical
consequences.  Sections~\ref{sec:v77-prefix}--\ref{sec:v77-coverage}
give the retained physical-action and whole-boundary proof chain.
Section~\ref{sec:v79-suspensions} proves the class-wide convex twist
result.  Section~\ref{sec:v78-mechanics} retains the full mass--potential
construction and its explicit formulas.  Sections~\ref{sec:v78-uniform}
and~\ref{sec:v79-geometric-robustness} give finite records, raw costs and
robust geometric output.

The periodic boundary-law inverse remains a separate companion.  It
retains the relative determinant normalization, positive quadratic
inverse, signed jets, actual smooth envelope and finite-preparation
consequences, including flat smooth perturbations.  Its determinant
estimate is a technical engine of that theorem, not a second independent
claim of significance for this article.  No input of that companion is
deleted or used as an unproved bridge between the two experiments.
